\printedpage{29}
% head: A. Grothendieck : Dérivateurs, Ch. XVII, Édition des Universités de Montpellier II et Paris VII
% foot: Version 9 février 2019
which are visibly adjoint to one another (since induced by a pair of adjoint functors % [sic] adjoint -> adjoints, accord avec « lesquels »
$u_!$, $u^*$), and moreover compatible with the localizers. So we can apply the adjoint functor lemma for the localized categories. To sum up:

\textbf{Corollary 3.5.} \textit{Suppose that $I$, $I'$ are artinian, and that we are in one of the following two cases:}
\begin{enumerate}
\item[a)] \textit{The spaces} $\Esp(I)$, $\Esp(I')$ \textit{are noetherian (e.g. $I$, $I'$ finite).}
\msp{page 58}
\item[b)] \textit{$C$ and $TC = C\cap W$ are stable under transfinite composition.}
\end{enumerate}
\textit{Suppose moreover that $u^*$ maps $\M(I'^\circ)_{\mathrm{c}}$ into $\M(I^\circ)_{\mathrm{c}}$, which is the case in particular if $u$ is locally cofiltering, i.e. the localized ordered sets $i'\backslash I$ are locally cofiltering. Then the functor $u_!^W$ constructed in (3.3.1) is left adjoint to $u_W^*$.}

\textbf{Remark 3.6.} When we have a pair of adjoint functors
\[
\begin{tikzcd}
\mathcal{C} \arrow[r, shift left, "u_!"] & \mathcal{C}' \arrow[l, shift left, "u^*"]
\end{tikzcd}
\ ,
\]
with the localizers $W$, $W'$ of $\mathcal{C}$, $\mathcal{C}'$, $u^*$ compatible with the localizers, and if moreover we have a full subcategory $\mathcal{C}_0$ of $\mathcal{C}$ such that
\begin{enumerate}
\item[a)] $u_!|\mathcal{C}_0$ transforms arrows of $W$ into arrows of $W'$, and
\item[b)] the inclusion $\mathcal{C}_0 \hookrightarrow \mathcal{C}$ induces an \textit{equivalence} $\mathcal{C}_0W_0^{-1} \xrightarrow{\sim} \mathcal{C}W^{-1}$ (where $W_0 = W\cap\Fl(\mathcal{C}_0)$),
\end{enumerate}
we find two canonical functors
\[
\begin{tikzcd}
\mathcal{C}W^{-1} \arrow[r, shift left, "u_!^W"] & \mathcal{C}'W'^{-1} \arrow[l, shift left, "u_W^*"]
\end{tikzcd}
\ ,
\]
where $u_!^W$ is defined only up to canonical isomorphism.

\msp{page 59}
Contrary to what \textsc{Anderson} claims, it is not true in general that these two % [sic] par -> pas, faute de frappe
functors are adjoint, without additional hypotheses. Here is a counterexample: we take $\mathcal{C} = \mathcal{C}'$, $u_! = u^* = \id_{\mathcal{C}}$, $W' = \operatorname{Is}(\mathcal{C})$, $W = \Fl(\mathcal{C})$, so $u_W^*$ is the canonical functor
\[
\underbrace{\Pi\,\mathcal{C}}_{\text{fundamental groupoid}} \qquad \longleftarrow \mathcal{C}\ .
\]
Suppose that $\mathcal{C}$ is 1-connected, i.e. $\Pi\,\mathcal{C}$ is equivalent to $e$. Let $e\in\Ob\mathcal{C}$ be such that $\operatorname{End}_{\mathcal{C}}(e)$ is the trivial group, and take $\mathcal{C}_0 = \{e\}$. Then the hypothesis on $\mathcal{C}_0$ is satisfied. But if $u_!^W$ were left adjoint to $u_W^*$, $e$ would be an initial object of $\mathcal{C}$. For example if $\mathcal{C}$ has a final object $e$ (which implies that $\mathcal{C}$ is 1-connected and $e$ rigid), this final object would have to be initial if \textsc{Anderson} were right with ``it is clear that \ldots''.

\msp{page 60}
Nevertheless:

\textbf{Theorem 3.7.} \textit{Let $(\M,W,C)$ be a category with cofibrations stable under finite $\varinjlim$, and $u : I \longrightarrow I'$ an increasing map between artinian ordered sets. We suppose that we are under one of the following hypotheses.}
\begin{enumerate}
\item[a)] $\Esp(I)$, $\Esp(I')$ \textit{are noetherian, e.g. $I$, $I'$ finite.}

\printedpage{30}
% head: A. Grothendieck : Dérivateurs, Ch. XVII, Édition des Universités de Montpellier II et Paris VII
% foot: Version 9 février 2019
\item[b)] \textit{$C$ and $TC \overset{\text{def}}{=} C\cap W$ are stable under transfinite composition, and $\M$ stable under small filtering $\varinjlim$.}
\end{enumerate}
\textit{Then the functor} $u_!^W : \Ho_W(I^\circ) \longrightarrow \Ho_W(I'^\circ)$ \textit{constructed in} (3.3.1) \textit{is left adjoint to the functor $u_W^*$ in the opposite direction.}

\textsc{Proof.} We use the well-known diagram
\[
\begin{tikzcd}[row sep=large, column sep=large]
 & I \arrow[d, "\partial_1"'] \arrow[dd, bend left=60, "w"] \\
I \arrow[r, shift left, "\partial_0\ (\text{open})"] \arrow[d, "u"'] \arrow[dr, phantom, "\text{cocart.}"] & I\times\boldsymbol{\Delta}^1 \arrow[l, dotted, shift left, "\pr_1"] \arrow[d] \\
I' \arrow[r, shift left, "\text{open imm.}"] & C(u) \arrow[l, dotted, shift left, "v"]
\end{tikzcd}
\]
to factor $u$ as
\[
u \;=\; vw\ ,
\]
where $w$ is a \textit{closed immersion,} and $v$ is obtained from $\pr_1 : I\times\boldsymbol{\Delta}^1 \longrightarrow I$ by push-out. It is a retraction of the mapping cone C(u) onto the open $I'$, whose closed complement is identified with $I$ via $w$.

\msp{page 61}
The gluing functor between $I'$, $I$
\[
H \;:\; I'^\circ\times I \;\longrightarrow\; \mathrm{Ens}
\]
is given by
\[
\underbrace{H(i',i)}_{=\,\Hom_{C(u)}(i',i)} \;=\; \Hom_{I'}(i',u(i))\ . \tag*{$(*)$}
\]
The retraction $v$ from $C(u)$ to $I'$ maps $i$ to $u(i)$, and an arrow $i'\longrightarrow i$ corresponding to $i'\longrightarrow u(i)$ in $I'$ via $(*)$, is transformed by $v$ into that same arrow $i'\longrightarrow u(i)$. Let us calculate, % [sic] transformé -> transformée, accord avec « une flèche »
for a given $i'_0$, the category $i'_0\backslash\Bigl(C(u),\ v : C(u)\longrightarrow I'\Bigr)$, of pairs of an object $\xi$ of $C(u)$ % [sic] des couple -> des couples, accord du pluriel
and an arrow $i'_0\longrightarrow v(\xi)$ in $I'$. We see that $i'_0\backslash C(u)$ admits the open $i'_0\backslash I'$, with the closed complement $i'_0\backslash I$,
\[
\begin{tikzcd}[column sep=large, row sep=huge]
I' \arrow[r, hook, "\text{open}"'] & C(u) \arrow[d] & I \arrow[l, hook', "\begin{array}{l}\text{closed}\\ \text{complement}\\ \text{of } I'\end{array}"'] \\
 & I'\ ,
\end{tikzcd}
\]
these two categories being glued by the bifunctor
\[
\begin{aligned}
(i'_0\backslash I')^\circ \times (i'_0\backslash I) &\quad\longrightarrow\quad \mathrm{Ens}\\
\Bigl((\alpha : i'_0\longrightarrow i'), (\beta : i'_0\longrightarrow u(i), i)\Bigr) &\quad\longmapsto\quad \Hom_{i'_0\backslash I'}(i',u(i))
\end{aligned}
\]
\printedpage{31}
% head: A. Grothendieck : Dérivateurs, Ch. XVII, Édition des Universités de Montpellier II et Paris VII
% foot: Version 9 février 2019
induced by $(*)$\,; i.e. the arrows from $(\alpha : i'_0\longrightarrow i') \in \Ob\,(i'_0\backslash I')$ to $(\beta : i'_0\longrightarrow u(i)) \in i'_0\backslash I$ are the arrows $i' \xrightarrow{\gamma} u(i)$ such that the triangle
\[
\begin{tikzcd}
 & i'_0 \arrow[dl, "\alpha"'] \arrow[dr, "\beta"] & \\
i' \arrow[rr, "\gamma"] & & u(i)
\end{tikzcd}
\]
commutes.

\msp{page 62}
If $i' = i'_0$, $\alpha = \id_{i'_0}$, then this commutativity amounts to $\gamma = \beta$, which means that there is one and only one arrow from $\alpha$ to $(i,\beta)$. This therefore shows that $(i'_0,\id_{i'_0})$ is an initial object of $i'_0\backslash C(u)$. We have therefore proved the lemma:

\textbf{Lemma 3.8.} \textit{Let $u : I\longrightarrow I'$ be an arrow in} (Cat)\textit{, let us factor it as}
\[
I \xrightarrow{\;w\;} C(u) \xrightarrow{\;v\;} I'\ ,
\]
\textit{where $C(u)$ is the mapping cone defined on page} 60\textit{, $w$ a closed immersion with complement $I'$, $v$ the canonical retraction of $C(u)$ onto $I'$. This being so, for every $i'_0\in\Ob I'$, the category $i'_0\backslash(C(u),v)$ admits an initial object – hence $v : C(u)\longrightarrow I'$ is (as thoroughly as can be) cofiltering.}

We can now prove 3.7. First of all, the hypotheses made on $I$, $I'$ carry over to $C(u)$ (artinian hypothesis, and possibly noetherian hypothesis on $\Esp(C(u))$). So
\msp{page 63}
we can apply the preceding results, in particular 3.5, to the arrows $v$, $w$. The preceding lemma indeed shows that $v$ satisfies the conditions of 3.5, so the functor $v_!^W$ constructed in 3.3.1 is left adjoint to $v_W^*$. Consider on the other hand the functor $w_!^W$ constructed in 3.3.1, I say that the same conclusion applies. This comes from the fact that $w$ is a closed immersion. Indeed, in this case it is immediate that
\[
w_! \;:\; \M(I^\circ) \;\longrightarrow\; \M(I'^\circ), \qquad\quad \text{where } I = C(u), % [sic] I = C(u) -> I' = C(u), w a pour but le mapping-cone C(u)
\]
is compatible with the localizers, since we have
\[
\left\{
\begin{array}{lcll}
w_!(X)(j) & \simeq & X(j) & \text{if } j\in I,\\[1ex]
w_!(X)(j) & \simeq & \emptyset_{\M} & \text{if } j\in J\backslash I,
\end{array}
\right.
\]
by the calculation of the fibers
\[
w_!(X)(j) \;\simeq\; \varinjlim_{i\in I/j} X(i)\ .
\]
So the definition of $w_!^W$ given in 3.3.1 amounts in fact to passing directly to the localized categories for $w_!$, without having to restrict ourselves to $\M(I^\circ)_{\mathrm{c}}$, and the adjoint functor lemma tells us that
\printedpage{32}
% head: A. Grothendieck : Dérivateurs, Ch. XVII, Édition des Universités de Montpellier II et Paris VII
% foot: Version 9 février 2019
\msp{page 64}
this functor is left adjoint to $w_W^*$.

Since we obviously have
\[
u_!^W \;\simeq\; v_!^W\circ w_!^W\ , \qquad\quad u_W^* \;\simeq\; w_W^*\circ v_W^*\ ,
\]
we conclude from this that $u_!^W$ is indeed left adjoint to $u_W^*$, q.e.d.

\textbf{Corollary 3.9.} \textit{Let $(\M,W,C)$ be a category with cofibrations. Then the prederivator $\mathbf{D}_W$, restricted to the category} $\mathrm{Diag}_0$ \textit{consisting of the categories $I$ that are} ordered, \textit{artinian and such that} $\Esp(I)$ \textit{is noetherian, satisfies the axioms} Der 1 (decomposition), Der 2 (localization), Der 3 (final object), Der 4 bis (existence of $u_!$), Der 5 bis (mode of calculation of the $u_!$), Der 6 (exactness of $W$-cocartesian squares). \textit{If $C$, $TC = C\cap W$ are stable under transfinite composition, we have the same result on the category} Diag \textit{consisting of the artinian ordered categories $I$, nothing more.}

\msp{page 65}
\textbf{3.10.} Thus, to have a derivator in the full sense of the term, either on $\mathrm{Diag}_0$, or on Diag, all that is missing is Der 4 (existence of the $u_*$), Der 5 (calculation of the $u_*$) and Der 6 (exactness of $W$-cartesian squares). We are going to obtain them in principle with the aid of a \textit{category with fibrations} structure $(\M,F,W)$ associated to $(\M,W)$, i.e. the kind of structure dual to the one introduced % [sic] associé -> associée, accord avec « structure »
on page 39. These two structures do not have to be related to each other in any way other than by the fact that they correspond to the same $W$. They are considerably easier to construct than Quillen structures $(W,C,F)$, where $C$, $F$ must be related in an extremely strict way, with the properties $C \longleftrightarrow \underbrace{TF}_{=\,F\cap W}\ ,\ \underbrace{TC}_{=\,C\cap W} \longleftrightarrow F$ (of which there is no mention in the axiomatics of K. S. \textsc{Brown}), and factorization for $(C,TF)$ and $(TC,F)$.

\msp{page 66}
But returning to the construction of the $u_!^W$, there remains the problem of constructing them for an arrow $u : I'\longrightarrow I$ \textit{whatsoever} in Cat. % [sic] u : I' -> I -> u : I -> I', u_!^W va de Ho_W(I°) dans Ho_W(I'°)

To tell the truth, the construction of the functor
\[
u_!^W \;:\; \Ho_W(I^\circ) \;\longrightarrow\; \Ho_W(I'^\circ) \tag*{(3.3.1)}
\]
is valid without supposing that $I'$ is an ordered category – it can be any category. This comes from the fact that the functor induced by $u_!$,
\[
u_{!\,\mathrm{cof}} \;:\; \M(I^\circ)_{\mathrm{c}} \;\longrightarrow\; \M(I'^\circ)
\]
is still compatible with the localizers, as we see by the calculation of the fibers
\[
u_!(X)(i) \;\simeq\; \varinjlim_{i\in I/i'} X(i)\ . % [sic] u_!(X)(i) -> u_!(X)(i'), la fibre est prise en i'
\]
Thus, we are reduced to the cases where $I' = e$, at the cost of replacing $I$ by $I/i'$, and noting that % [sic] aux cas -> au cas, singulier
$(I/i'\longrightarrow I$ being étale) the inverse image of $X$ on $I/i'$ is cofibrant (cf. 3.2.a). % [sic] cofibrant -> cofibrante, accord avec « l'image inverse »

\msp{page 67}
But there arises again the perplexity whether the functor (3.3.1) thus constructed is indeed left adjoint to $u_W^*$. We can think of using once again the same factorization as $u = vw$, by
\printedpage{33}
% head: A. Grothendieck : Dérivateurs, Ch. XVII, Édition des Universités de Montpellier II et Paris VII
% foot: Version 9 février 2019
using the mapping cone. The difficulty, however, is that now $J = C(u)$ is no longer an ordered category, so we do not have a ready-made definition of $C_{J^\circ}\subset\Fl(\M(J^\circ))$, giving rise to a subcategory $\M(J^\circ)_{\mathrm{c}}$ etc. In fact, it would seem that we are practically paralyzed, so long as we do not manage to construct on $\M(J^\circ)$ also a category with cofibrations structure, behaving well with respect to the functors $v_!$. The % [sic] au foncteurs -> aux foncteurs, accord du pluriel
``proof'' by \textsc{Anderson} of the existence of $u_!$ for \textit{every} arrow in Cat, is a total flop (even for $I\longrightarrow e$, $I$ a finite ordered set – but there I was able to set things right), for he reduces to the case where $I$ is ordered and locally finite, and he
\msp{page 68}
assumes without further ado that \textit{in this case} $u_!^W$ is left adjoint to $u_W^*$, and \textit{moreover,} that if $u$ is $\mathbf{W}_0$-connected, hence $u^* : \M(I'^\circ)\longrightarrow\M(I^\circ)$ fully faithful, that the same is true of $u_W^*$ obtained by passing to the localized categories. Now even if we grant that $u_!^W$ is left adjoint to
$u_*^W$, for lack of having an adjunction morphism close at hand ``down below'' (in $\Ho_W(I'^\circ)$), % [sic] u_*^W -> u_W^*, notation employée partout ailleurs
obtained from the adjunction morphism in $\M(I^\circ)$, it is not at all clear that $u_W^*$ is fully faithful, i.e. $u_!^W$ a localization functor. Counterexample 3.6 clearly shows that, even if we have adjunction (i.e. if $e$ is an initial object of $\mathcal{C} = \mathcal{C}'$), the functor $u^* = \id_{\mathcal{C}}$ may well be fully faithful (and even an equivalence of categories), the localized functor has no reason to be so: this would prove that every object of Cat having an initial object % [sic] pouverait -> prouverait, faute de frappe
is equivalent to $e$\,!

\msp{page 69}
\textbf{3.11.} Let us say that we have a \textit{strict category with cofibrations,} if in the factorization axiom $\mathrm{Cof}_3$ (cf. p. 39), the factorization can be taken functorially in $f$. It amounts to $\pm$ the same thing to say that the factorization axiom remains valid, when we think of a category $\M(I)$ for $I$ in Cat, defining $W(I)$ and $C(I)$ there (not to be confused with $C_I$) fiber by fiber. This amounts to the same thing at least if $\M$ is equivalent to a small category. Otherwise, we must assume functoriality for every \textit{small} full subcategory of $\underline{\mathrm{Fl}}(\M)$. It is better to define the strict category with cofibrations by this stability property.

We then see that the prederivator $\mathbf{D} = \mathbf{D}_{\M}$ on Cat is a derivator \textit{relative to} $\mathcal{D} = \mathrm{Diag}_0$ defined above (artinian ordered categories $I$, such that $\Esp(I)$ is noetherian), or even relative to $\mathcal{D} = \mathrm{Diag}$ (artinian ordered categories $I$, nothing more) when $C$ and $TC$ are stable under transfinite composition. By this I mean that $u_!$ exists when $u$ is an arrow of $\mathcal{D}$,
\msp{page 70}
and that this remains valid when we pass to an induced derivator $\mathbf{D}_Z$ ($Z$ arbitrary in Cat)\,; in other words, that $v_!$ exists when $v$ is of the form $\id_Z\times u$, where $u : I\longrightarrow I'$ is an % [sic] d'autre termes -> d'autres termes, accord du pluriel
arrow of $\mathcal{D}$, and that the calculation is done \textit{over} $I'$ according to the rule Der 5 bis. This therefore means that the base change arrow for the $f_!$ is an isomorphism in the case % [sic] un isomorphismes -> un isomorphisme, accord du singulier
that follows:
\[
\begin{tikzcd}[row sep=huge, column sep=huge]
I\times Z \arrow[d, "u\times\id_Z"] & (i'\backslash I)\times Z \arrow[l] \arrow[d, "\pr_2"] \\
I'\times Z & \{i'\}\times Z \arrow[l]
\end{tikzcd}
\ ,
\]
where $i'$ is an arbitrary object of $I'$.